\section{Motivation}
The primary motivation for this work is the need to reduce communication barriers faced by ISL users in everyday interactions. In academic institutions, hospitals, offices, and public service centers, communication often depends on interpreter availability, which may not always be practical. A robust automated recognition system can act as an assistive bridge, enabling faster and more inclusive interaction between signers and non-signers.

From a technical perspective, ISL recognition is a meaningful and challenging problem at the intersection of computer vision, sequence modeling, and human-centered AI. It requires the model to understand both fine-grained visual structure and temporal progression under real-world variability. Addressing this challenge contributes not only to assistive communication technology but also to broader research on efficient multimodal understanding systems.

This thesis is therefore motivated by a combined social and engineering objective: to build a system that is both practically useful and technically sound, with a clear pathway toward future real-world deployment.
